\documentclass[11pt,letterpaper]{article}
\usepackage{fontspec}
\usepackage{tocloft}
\usepackage{multicol}
\usepackage{nicefrac}
\usepackage[left=1.4in,right=1.5in,top=1in,bottom=1in]{geometry}
\setmainfont[Scale=1.4,AutoFakeBold=1.5,AutoFakeSlant=0.3]{Doves Type}

\widowpenalty=10000
\clubpenalty=10000
\interlinepenalty=500

\title{Shallot-Dijon Cream Sauce}
\author{}
\date{}

\begin{document}

\maketitle
\thispagestyle{empty}

\begin{quote}
\textit{Shallots are sweated in butter until deeply softened, then deglazed with dry white wine and finished with heavy cream and Dijon mustard. The result is a pourable, coating sauce with gentle tang and herb finish---ideal alongside chicken pasties, roasted poultry, or steamed vegetables. Yields about 2~cups.}
\end{quote}

\section*{Ingredients}

\setlength{\columnsep}{20pt}
\begin{multicols}{2}
\noindent
    Shallots, medium \dotfill 3 \\
    Unsalted butter \dotfill 3 Tbsp. \\
    Dry white wine \dotfill \nicefrac{1}{3} cup \\
    Heavy cream \dotfill 1\nicefrac{1}{3} cups \\
    \columnbreak
    Dijon mustard \dotfill 3--4 tsp. \\
    Kosher salt \dotfill to taste \\
    Black pepper, freshly ground \dotfill to taste \\
    Fresh tarragon or chives, minced \dotfill 1--2 Tbsp.
\end{multicols}

\section*{Directions}

\noindent
Mince \textbf{shallots} and have ready --- Measure \textbf{wine} and \textbf{cream} --- Have \textbf{Dijon}, \textbf{salt}, and \textbf{pepper} nearby --- Mince \textbf{tarragon} or \textbf{chives} and set aside in \textit{Small Bowl~\#1}

\begin{enumerate}
    \item In a medium saucepan or skillet over \textit{medium-low heat}, melt 3~Tbsp. \textbf{butter}. Add minced \textbf{shallots} and cook, stirring occasionally, until deeply softened and translucent with no raw bite, \textit{8--10~minutes}. Do not brown; reduce heat if shallots begin to color. They should be tender and sweet.

    \item Add \nicefrac{1}{3}~cup \textbf{dry white wine} and increase heat to \textit{medium}. Scrape any bits from the bottom and simmer until wine is reduced by half, \textit{2--3~minutes}. The pan should smell winey but not sharp.

    \item Add 1\nicefrac{1}{3}~cups \textbf{heavy cream} and bring to a gentle simmer over \textit{medium} or \textit{medium-low heat}. Cook, stirring occasionally, until sauce thickens enough to coat the back of a spoon and leaves a clear trail when you run a finger through it, \textit{3--5~minutes}. Do not boil vigorously or the cream may separate.

    \item Remove from heat. Whisk in 3--4~tsp. \textbf{Dijon mustard} (to taste). Season with \textbf{kosher salt} and \textbf{black pepper}. Taste and adjust---sauce should be well-seasoned, with a clear but not harsh mustard note.

    \item Stir in minced \textbf{tarragon} or \textbf{chives} (\textit{Small Bowl~\#1}). If sauce is thicker than desired (e.g., not pourable for drizzling), thin with 1--2~Tbsp. chicken stock or milk and reheat gently. Serve warm.
\end{enumerate}

\newpage

\section*{Equipment Required}
\begin{itemize}
    \item Medium saucepan or skillet (2--3~qt.)
    \item Whisk
    \item Measuring cups and spoons
    \item Cutting board and sharp knife
    \item Small bowl for minced herbs
    \item Rubber spatula
\end{itemize}

\section*{Hints and Notes}

\subsection*{Yield}
\begin{itemize}
    \item Makes about 2~cups sauce
    \item Enough for 8--10 pasties with generous drizzling or dipping, or 4--6 portions as a side sauce for poultry or vegetables
\end{itemize}

\subsection*{Mise en Place}
\begin{itemize}
    \item Mince \textbf{shallots} and have ready before heating the pan
    \item Measure \textbf{wine} and \textbf{cream} and have nearby for time-sensitive steps
    \item Mince \textbf{tarragon} or \textbf{chives} and set aside in \textit{Small Bowl~\#1} for the finish
\end{itemize}

\subsection*{Ingredient Tips}
\begin{itemize}
    \item Use \textbf{dry white wine} you would drink---Pinot Grigio, Sauvignon Blanc, or unoaked Chardonnay work well; avoid ``cooking wine'' with added salt
    \item \textbf{Crème fraîche} can replace part of the \textbf{heavy cream} for extra tang; adjust consistency with stock or milk if needed
    \item \textbf{Fresh tarragon} pairs classically with cream and mustard; \textbf{chives} offer a milder, oniony finish
\end{itemize}

\subsection*{Preparation Tips}
\begin{itemize}
    \item Sweating \textbf{shallots} slowly (without browning) is key---they should be soft and sweet, not caramelized or bitter
    \item Reducing the \textbf{wine} by half removes raw alcohol and concentrates flavor without harshness
    \item If the sauce becomes too thick after cooling, thin with chicken stock or milk and reheat gently; do not boil after adding herbs
\end{itemize}

\subsection*{Make Ahead \& Storage}
\begin{itemize}
    \item Refrigerate in an airtight container for \textit{up to 5~days}
    \item Reheat gently over \textit{low heat}, stirring; thin with a splash of cream or stock if needed
    \item Fresh herbs are best added just before serving; if making ahead, add them when reheating
\end{itemize}

\subsection*{Serving Suggestions}
\begin{itemize}
    \item Serve alongside Chicken, Mushroom \& Spinach Pasties for dipping or drizzling
    \item Spoon over roasted or pan-seared chicken, turkey cutlets, or pork tenderloin
    \item Drizzle over steamed asparagus, green beans, or broccoli
    \item Use as a warm dressing for blanched vegetables or as a base for a quick pasta sauce (thin slightly and toss with pasta and peas or ham)
\end{itemize}

\subsection*{Wine Pairings}
\begin{itemize}
    \item \textbf{With the sauce and pasties:} Choose wines that match the richness and slight tang---unoaked or lightly oaked \textbf{Chardonnay}, \textbf{Sauvignon Blanc} (dry), \textbf{Pinot Gris} / \textbf{Pinot Grigio}, or \textbf{Albariño}. Dry \textbf{Riesling} (Kabinett or dry Alsace style) also works.
    \item Avoid heavily oaked or very sweet wines; the sauce's mustard and cream are best with crisp, clean whites or light-bodied options.
\end{itemize}

\end{document}
